% sec-carry.tex -- the shared carry-chain device for the
% kp+1 palindrome developments (k in {2,3,4,5,6}).
% Standalone fragment, \input by the master file.
% The theorem environments are declared here (first input
% in document order); the later section files guard their
% own declarations with \@ifundefined so that every
% fragment also compiles standalone.

\theoremstyle{plain}
\newtheorem{theorem}{Theorem}
\newtheorem{lemma}[theorem]{Lemma}
\newtheorem{corollary}[theorem]{Corollary}
\newtheorem{question}[theorem]{Question}
\theoremstyle{remark}
\newtheorem{remark}[theorem]{Remark}
\providecommand{\IsPal}{\mathrm{IsPal}}

\section{The shared carry-chain device}
\label{sec:carry}

Fix $k \in \{2,3,4,5,6\}$. Let $p$ be a decimal
palindrome of $m \ge 1$ digits, and assume $q = kp+1$ is
again a palindrome, of $m'$ digits. Write
$a_j = \lfloor p/10^j \rfloor \bmod 10$ for the $j$-th
digit of $p$ (units $= j = 0$) and define the carry into
position $j$ of the computation of $kp+1$ (the usual
base-$10$ addition carry, cf.~\cite{holte}) by the closed
form
\begin{gather*}
 c_j \;=\; \bigl\lfloor\bigl(k\,(p \bmod 10^j) + 1\bigr)\,
 \big/\,10^j\bigr\rfloor ,
 \\
 \text{so that}\qquad
 b_j = (k a_j + c_j) \bmod 10 ,
 \qquad c_{j+1} = \bigl\lfloor\bigl(k a_j + c_j\bigr)\big/10
 \bigr\rfloor ,
\end{gather*}
where $b_j$ is the $j$-th digit of $q$. Then $c_0 = 1$
(the $+1$ is folded into the chain), $0 \le c_j \le k-1$,
and since $kp+1 < 10^{m+1}$ the length satisfies
$m' \in \{m, m+1\}$, with $m' = m$ iff $c_m = 0$.
Palindromicity of $q$ reads $b_j = b_{m'-1-j}$. We write
$R_m = (10^m-1)/9$ for the $m$-digit repunit and
$d\,R_m$ for the repdigit $dd\dots d$. We write
$\IsPal(p,m)$ for the predicate that $p$ is a positive
$m$-digit integer ($1 \le p \le 9$ for $m = 1$, else
$10^{m-1} \le p < 10^m$) whose base-$10$ digits satisfy
$a_j = a_{m-1-j}$ for all $0 \le j < m$.

\begin{lemma}[same length: symmetric $2$-periodic chain]
\label{lem:same}
If $m' = m$ then, for every $k \in \{2,3,4,5,6\}$:
\emph{(i)} $c_j = c_{m-1-j}$ for all $j < m$;
\emph{(ii)} $c_{j+1} = c_{j-1}$ for $1 \le j \le m-2$,
hence $c_{2t} = 1$ for $2t \le m-1$ and
$c_{2t+1} = c_1$ for $2t+1 \le m-1$;
\emph{(iii)} $m$ is odd and $c_1 = 0$;
\emph{(iv)} the digits obey the alternating classes
$a_j \le \lfloor 8/k \rfloor$ at even positions and
$\lceil 10/k \rceil \le a_j \le \lfloor 19/k \rfloor$ at
odd positions; in particular
$a_0 \in \{1,\dots,\lfloor 8/k \rfloor\}$.
\end{lemma}

\begin{proof}
(i) $b_j = b_{m-1-j}$ reads
$(k a_j + c_j) \equiv (k a_{m-1-j} + c_{m-1-j})
\pmod{10}$; the input symmetry $a_{m-1-j} = a_j$ cancels
$k a_j$ exactly, and the residues
$c_j, c_{m-1-j} \le k-1 < 10$ coincide. (ii) Apply the
recurrence at $j$ and at the mirror index $m-1-j$, then
rewrite through (i) and the same digit symmetry. (iii) If
$m$ is even, $m-1$ is odd: (i) at $j = 0$ gives
$c_{m-1} = c_0 = 1$ while (ii) gives $c_{m-1} = c_1$, so
the top recurrence reads
$c_m = \lfloor(k a_0 + c_{m-1})/10\rfloor
= \lfloor(k a_0 + 1)/10\rfloor = c_1 = 1$,
contradicting $c_m = 0$. Hence $m$ is odd, $m-1$ is
even, $c_{m-1} = 1$ by (ii), and
$0 = c_m = \lfloor(k a_0 + 1)/10\rfloor$ forces
$k a_0 \le 8$ and $c_1 = 0$. (iv) At even $j$ one has
$c_j = 1$ and $c_{j+1} = 0$ (odd index, or $j+1 = m$), so
$\lfloor(k a_j + 1)/10\rfloor = 0$, i.e.\ $k a_j \le 8$.
At odd $j$, $c_j = c_1 = 0$ and $c_{j+1} = 1$ (even
index), so $10 \le k a_j \le 19$.
\end{proof}

Table~\ref{tab:classes} lists the alternating classes of
Lemma~\ref{lem:same}(iv) for each multiplier.

\begin{table}[ht]
\centering
\small
\begin{tabular}{c|c|c|c|c}
\hline
$k$ & entry & even positions & odd positions &
$a_0$ (same length) \\
\hline
2 & A069881 & $\{0,1,2,3,4\}$ & $\{5,6,7,8,9\}$ &
$\{1,2,3,4\}$ \\
3 & A083829 & $\{0,1,2\}$ & $\{4,5,6\}$ & $\{1,2\}$ \\
4 & A083831 & $\{0,1,2\}$ & $\{3,4\}$ & $\{1,2\}$ \\
5 & A083833 & $\{0,1\}$ & $\{2,3\}$ & $\{1\}$ \\
6 & A083835 & $\{0,1\}$ & $\{2,3\}$ & $\{1\}$ \\
\hline
\end{tabular}
\caption{Digit classes forced on same-length members, by
multiplier. Even-length members are excluded from this
table: they all live in the overflow branch
(Lemma~\ref{lem:overflow}).}
\label{tab:classes}
\end{table}

\begin{lemma}[sufficiency of the alternating classes]
\label{lem:altsuff}
Let $k \in \{2,3,4,5,6\}$, let $m$ be odd, and let $p$ be
an $m$-digit palindrome whose digits belong to the
alternating classes of Table~\ref{tab:classes}: $a_0 \ge
1$, $k a_j \le 8$ at even positions and
$10 \le k a_j \le 19$ at odd positions. Then $kp+1$ is an
$m$-digit palindrome.
\end{lemma}

\begin{proof}
The chain closes by induction: $c_0 = 1$; at an even
position $k a_{2t} \le 8$ gives $b_{2t} = k a_{2t} + 1 \le
9$ and $c_{2t+1} = 0$; at an odd position
$10 \le k a_{2t+1} \le 19$ gives $b_{2t+1} = k a_{2t+1} -
10$ and $c_{2t+2} = 1$. Hence $c_{2t} = 1$ and
$c_{2t+1} = 0$ throughout; as $m-1$ is even the top carry
is $c_m = 0$, so $q = kp+1$ keeps length $m$. Since $j$
and $m-1-j$ have the same parity and the output digit
depends only on the parity of the position
($b_{2t} = k a_{2t} + 1$, $b_{2t+1} = k a_{2t+1} - 10$),
$a_{m-1-j} = a_j$ gives $b_{m-1-j} = b_j$.
\end{proof}

\begin{lemma}[overflow: the frozen chain]
\label{lem:overflow}
If $m' = m+1$ the head equation
\begin{equation}
 (k a_0 + 1) \bmod 10 \;=\; c_m ,
 \qquad c_m = \bigl\lfloor\bigl(k a_0 + c_{m-1}\bigr)\big/10
 \bigr\rfloor
 \label{eq:head}
\end{equation}
(the first digit $b_0$ equals the new leading digit
$b_m = c_m$) has, for each $k \in \{2,3,4,5,6\}$, a single
solution: $a_0 = d_k$ and $c_m = v_k$, where $(d_k, v_k)$
is the unique digit solution of $k d + 1 = 11 v$, namely
\[
 (k;\,d_k;\,v_k) \;=\;
 (2;5;1),\; (3;7;2),\; (4;8;3),\; (5;2;1),\; (6;9;5) .
\]
The whole chain freezes: $c_j = v_k$ for $1 \le j \le m$
and every digit of $p$ equals $d_k$. Hence $p = d_k\,R_m$
is a repdigit, and $q = kp+1$ is the repunit $R_{m+1}$
for $k \in \{2,5\}$, respectively the framed repdigit
$v_k\,w_k^{\,m-1}\,v_k$ with
$w_k = (k d_k + v_k) \bmod 10$ for $k \in \{3,4,6\}$,
i.e.\ $2\,3^{\,m-1}\,2$, $3\,5^{\,m-1}\,3$ and
$5\,9^{\,m-1}\,5$.
\end{lemma}

\begin{proof}
\emph{Scheme I} ($k \in \{2,3,4,6\}$). The head equation
\eqref{eq:head} is solved by a finite check on
$a_0 \in \{1,\dots,9\}$ against $1 \le c_m \le k-1$ and
$c_m = \lfloor(k a_0 + c_{m-1})/10\rfloor$: only
$a_0 = d_k$ survives, giving $c_m = v_k$ and $c_1 = v_k$.
The propagation at position $j+1$ uses the mirror couple
\emph{(A)}
$(k a_{j+1} + c_{j+1}) \equiv (k a_j + c_{m-j-1})
\pmod{10}$ (the output symmetry folded by
$a_{m-j-1} = a_j$) and \emph{(B)}
$c_{m-j-1} = \lfloor(k a_{j+1} + c_{m-j-2})/10\rfloor$
(the mirrored recurrence). For $k \in \{2,4,6\}$ the pair
\emph{(A)}+\emph{(B)} has the single fixed point
$(a_{j+1}, c_{j+2}) = (d_k, v_k)$: the chain is rigid in
one step (for $k = 2$ the parity of $2a'+1$ forces the
mirror carry to $1$, then $a' \in \{0,5\}$ meets
$a' \ge 5$; for $k = 4$ and $k = 6$ the joint system is
solved digit by digit over the ten possibilities of
$a'$). For $k = 3$ the state $(a',c') = (3,1)$ is a
spurious fixed point of \emph{(A)}+\emph{(B)}; it is
killed one position later by the same equations shifted:
the shifted mirror congruence forces the next digit to
be $6$ (so $3 \cdot 6 = 18$), and the shifted mirrored
recurrence would then require the carry
$0 = \lfloor(18 + c)/10\rfloor$, impossible, so the induction runs
over $j \le m-2$ and the last digit is closed by
palindromy $a_{m-1} = a_0$.
\emph{Scheme II} ($k = 5$). The multiplier $5$ separates
carries mod $10$: $5a \bmod 10 \in \{0,5\}$ leaves
$c < 5$ readable in $b_j$, so the mirror congruence
yields directly $a_j \equiv a_{j-1} \pmod 2$ \emph{and}
$c_j = c_{m-j}$ ($1 \le j \le m-1$). A $2$-cycle
$g(u,a) = w$, $g(w,a) = u$ with $u \ne w$ is impossible
($11\,|u-w| \le 9$), hence the chain is constant
$c_j = v$ without induction; \eqref{eq:head} forces
$v = 1$ and $a_0 = 2$, and
$\lfloor(5 a_j + 1)/10\rfloor = 1$ with the common parity
gives $a_j = 2$ throughout. In all cases $b_0 = b_m = v_k$
and $b_j = (k d_k + v_k) \bmod 10$ inside, which is the
displayed image form.
\end{proof}

\begin{remark}[what the factorization actually shares]
\label{rem:honest}
The same-length half (Lemma~\ref{lem:same}) is literally
the same proof in all five certified developments; the
overflow half factors the \emph{statement} (frozen
constant chain, repdigit member) but not the proof:
$k \in \{2,4,6\}$ is rigid in one step, $k = 3$ needs the
one-step dead-end lookahead, and $k = 5$ goes through the
preserved mirror symmetry and the no-$2$-cycle argument,
a route closed to the other multipliers where
$k a \bmod 10$ does not separate the carry. In the kernel
development the machinery is likewise duplicated per
multiplier: there is no parametrized general lemma, the
same carry-chain argument is formalized for each $k$
(Section~\ref{sec:lean}).
\end{remark}

\begin{remark}[the multipliers $k = 1$ and
$k \in \{7,8,9\}$]
\label{rem:k1}
For $k = 1$ the class is finite: $S_1 = \{1,\dots,8\}$.
An $m$-digit palindrome with $m \ge 2$ and $a_0 \le 8$
would give $p+1$ the same leading digit $a_0$ but last
digit $a_0+1$, never palindromic; $a_0 = 9$ makes $p+1$
end in $0$, never a palindrome either. This elementary
case needs no machine. For $k \in \{7,8,9\}$ the head
key $kd+1 = 11v$ still has solutions -- $(d,v) = (3,2)$,
$(4,3)$, $(6,5)$ -- predicting repdigit overflow families
$33\ldots3 \mapsto 2\,3^{\,m-1}\,2$,
$44\ldots4 \mapsto 3\,5^{\,m-1}\,3$ and
$66\ldots6 \mapsto 5\,9^{\,m-1}\,5$; the displayed terms
of A083837 agree ($a(7) = a(8) = 1$, $a(9) = 6$, and
$7\cdot33+1 = 232$, $8\cdot44+1 = 353$,
$9\cdot66+1 = 595$ are palindromes). The full
classification of those entries is outside the certified
perimeter of this paper.
\end{remark}
